% VI35-DETAILED-PROBLEM-15
\begin{problem}[VI.35.P15: Irrelevant torsion sheafifies to zero]\label{prob:vi35-15}
Let \(M\) be graded and suppose \(S_+^NM=0\) for some \(N\). Prove \(\widetilde M=0\). Give a concrete example.

\textbf{Provenance.} Canonical VI/35 extension.
\end{problem}

\ifdefined\FullProblemDossiers
\begin{solution}
For every homogeneous \(f\in S_+\),
\[
f^N\in S_+^N,
\]
so
\[
f^NM=0.
\]
After localizing at \(f\), the element \(f^N\) becomes a unit. Hence
\[
M_f=0,
\]
and therefore
\[
M_{(f)}=(M_f)_0=0.
\]
All standard-chart restrictions of \(\widetilde M\) vanish. Since standard opens cover \(\Proj S\),
\[
\boxed{\widetilde M=0}.
\]

For a concrete example take
\[
S=k[x,y],
\qquad
M=S/S_+=S/(x,y)
\]
with its induced grading. Then
\[
S_+M=0,
\]
so \(\widetilde M=0\) on \(\Proj S\), even though \(M\neq0\) as a graded module. Thus the projective tilde functor is
not faithful on the entire category of graded modules: it discards module information supported only at the
irrelevant cone vertex.
\end{solution}
\fi
